% Bagean: sajarah
% Bab: biografi
% Pérangan: julia-robinson
\documentclass[../../../include/open-logic-section]{subfiles}

\begin{document}

\olfileid{his}{bio}{rob}

\olsection{Julia Robinson}

\olphoto{robinson-julia}{Julia Robinson}

Julia Bowman Robinson iku matématikawan Amérika.
Dhèwèké mligi kawentar amarga pakaryané ngenani
masalah putusan, lan sing paling misuwur amarga
sumbangané kanggo ngrampungaké masalah kaping
sepuluh Hilbert. Robinson lair ing St.~Louis,
Missouri, tanggal 8 Desember 1919. Robinson
kèlingan yèn wiwit cilik wis kasengsem marang
wilangan \citep[4]{Reid1986}. Nalika umur
sangang taun dhèwèké kena demam scarlet lan
bola-bali ngalami demam rematik. Iki meksa
dhèwèké ngentèkaké akèh wektu ing amben,
saéngga ketinggalan sekolah. Senajan bisa
nguber ketinggalan kanthi pitulungan tutor
pribadi, akibat fisik saka penyakité
tetep mengaruhi uripé nganti suwé.

Senajan ngalami kasulitan nalika cilik,
Robinson lulus sekolah menengah kanthi
pirang-pirang penghargaan matématika lan
ilmu pengetahuan. Dhèwèké miwiti kuliah
ing San Diego State College, banjur pindhah
menyang University of California, Berkeley,
nalika taun pungkasan studi sarjanané.
Ing kono dhèwèké kapengaruh déning
matématikawan Raphael Robinson. Kekaroné
dadi kanca cedhak lan omah-omah ing taun
1941. Minangka pasangan saka anggota
staf pengajar, Robinson dilarang mulang
ing dhépartemèn matématika Berkeley.
Senajan tetep ngrungokaké kuliah-kuliah
matématika, dhèwèké ngarep-arep bisa
ninggalaké universitas lan miwiti
kulawarga. Nanging ora suwé sawisé
omah-omah, Robinson kena pneumonia.
Dhèwèké dikandhani yèn akèh jaringan
parut ngumpul ing jantungé amarga
demam rematik nalika cilik. Amarga
parut iku abot, dhokter prédhiksi
dhèwèké ora bakal urip ngliwati
umur patang puluh taun lan menehi
pitutur supaya ora duwé anak
\citep[13]{Reid1986}.

Robinson ngalami depresi nganti suwé,
nanging pungkasané mutusaké nerusaké
sinau matématika. Dhèwèké bali
menyang Berkeley lan ngrampungaké
gelar doktor ing taun 1948 ing
sangisoré bimbingan Alfred Tarski.
Tarski wis nuduhaké yèn téori
orde kapisan kanggo wilangan real
bisa diputusaké, lan saka karya
G\"odel bisa disimpulaké yèn
téori orde kapisan kanggo
wilangan asli ora bisa
diputusaké. Pitakonan gedhé
kang durung kajawab yaiku
apa téori orde kapisan
kanggo wilangan rasional
bisa diputusaké. Ing
disertasiné \citeyearpar{Robinson1949},
Robinson mbuktèkaké yèn
ora bisa.

Amarga kasengsem marang masalah
putusan, Robinson banjur nyoba
ngrampungaké masalah kaping sepuluh
Hilbert. Masalah iki kalebu
dhaptar 23 masalah matématika
kondhang kang diajokaké déning
David Hilbert ing taun 1900.
Masalah kaping sepuluh takon
apa ana algoritma kang bisa
mangsuli sajroning wektu winates
apa sawijining persamaan
polinomial kanthi koefisièn
wilangan wutuh, kayata
$3x^2 - 2y +3 = 0$,
nduwèni solusi ing wilangan
wutuh. Pitakonan kaya iki
diarani \emph{masalah Diofantin}.
Sawisé sawetara kasil wiwitan,
Robinson kerja bareng karo
Martin Davis lan Hilary Putnam,
kang uga nggarap masalah iku.
Wong telu iku kasil nuduhaké
yèn masalah Diofantin
èksponènsial (ing ngendi
peubah uga bisa dadi pangkat)
ora bisa diputusaké, lan
nuduhaké yèn sawijining
konjektur (banjur diarani
``J.R.'') ndadèkaké
masalah kaping sepuluh
Hilbert ora bisa diputusaké
\citep{DavisPutnamRobinson1961}.
Robinson terus nggarap
masalah iki sajroning
dasawarsa 1960-an. Ing taun
1970, matématikawan enom
saka Rusia Yuri Matijasevich
pungkasané mbuktèkaké
hipotèsis J.R. Asil
gabungan iku saiki diarani
teorema
Matijasevich--Robinson--Davis--Putnam,
utawa cekaké teorema MRDP.
Matijasevich lan Robinson
dadi kanca lan kerja bareng
nulis sawetara makalah.
Ing sawijining layang
marang Matijasevich,
Robinson tau nulis yèn
``satemené aku bungah
banget amarga kanthi
kerja bareng (senajan
pisah éwonan mil) kita
cetha maju luwih adoh
tinimbang yèn saben
wong nyambut gawé
dhéwé-dhéwé''
\citep[45]{Matijasevich1992}.

Robinson iku wanita kapisan
kang dadi présidhèn
American Mathematical
Society, lan wanita
kapisan kang kapilih
ing sèksi matématika
National Academy
of Sciences. Dhèwèké
séda tanggal 30 Juli
1985 nalika umur 65
taun, sawisé
didiagnosis ngalami
leukemia.

\begin{reading}
Makalah-makalah matématika
Robinson kasedhiya ing
\textit{Collected
  Works} \citep{Robinson1996},
kang uga ngemot cetakan
ulang memoir biografiné
saka National Academy of
Sciences
\citep{Feferman1994}.
Mbakyuné Robinson,
Constance Reid, nerbitaké
``Autobiography of Julia,''
adhedhasar wawancara
\citep{Reid1986},
uga sawijining memoir
jangkep \citep{Reid1996}.
Dhokumèntèr cekak
ngenani Robinson lan
masalah kaping sepuluh
Hilbert disutradarani
George Csicsery
\citep{Csicsery2016}.
Kanggo memoir ringkes
ngenani kerja bareng
Yuri Matijasevich karo
Robinson lan pangaruh
Robinson marang pakaryané,
delengen
\citep{Matijasevich1992}.
\end{reading}

\end{document}
